% GENERATED FILE. Edit canonical lessons, then run npm run generate:books.
% canonical-source-hash: fnv1a64:52ee4d0649b94bb0
% canonical-lessons: ML-C206-kusrtiyulla, ML-C206-anusaranayulla, ML-C206-svartthamaya, ML-C206-sthiramaya, ML-C206-orra

\chapter{Naughty, Obedient, Selfish, Permanent, Single}
\label{ch:ml-naughty-obedient-selfish-permanent-single}
\hlchaptermodality{206}

\begin{chapteropening}
\textbf{By the end of this chapter:} \emph{I can say naughty, obedient, selfish, permanent and single.}

Complete the last lesson of chapter 206: I can say naughty, obedient, selfish, permanent and single.
\end{chapteropening}

\section[\texorpdfstring{\ml{കുസൃതിയുള്ള} (kusṛtiyuḷḷa)}{കുസൃതിയുള്ള (kusṛtiyuḷḷa)}]{\ml{കുസൃതിയുള്ള} (kusṛtiyuḷḷa) — naughty, mischievous}

\label{lesson:ML-C206-kusrtiyulla}



\begin{quote}
Before the new one: say the Malayalam for strange, then the Malayalam for modern.
\end{quote}

\subsection*{You'll want to know: \ml{കുസൃതിയുള്ള}}
\textbf{\ml{കുസൃതിയുള്ള}} — \emph{kusṛtiyuḷḷa} — \textquotedblleft{}naughty, mischievous\textquotedblright{}.

\begin{cousinweb}
\textbf{\ml{കുസൃതിയുള്ള}} (\emph{kusṛtiyuḷḷa}) joins Sanskrit \emph{kusṛti}, \textquotedblleft{}trickery, a crooked way\textquotedblright{}, to native \emph{-uḷḷa}, \textquotedblleft{}that has\textquotedblright{} (DEDR 697).
\end{cousinweb}

\begin{grammarlens}[title={What you've built}]
One more word for this chapter.
\end{grammarlens}

\subsection*{Your turn}
\begin{itemize}
\raggedright
  \item \emph{Say it:} \emph{kusṛtiyuḷḷa}
  \item \emph{Say it:} \emph{kusṛtiyuḷḷa}, once more
  \item \emph{Say it:} say \emph{kusṛtiyuḷḷa}
  \item \emph{Recall:} say the Malayalam for strange, then the Malayalam for modern, then say \emph{kusṛtiyuḷḷa} again
\end{itemize}

\begin{tcolorbox}[breakable,colback=teal!4,colframe=teal!35!black,title={Before you move on}]
What does \ml{കുസൃതിയുള്ള} mean? (\textquotedblleft{}Naughty, mischievous\textquotedblright{}.) Say it once more.
\end{tcolorbox}
\section[\texorpdfstring{\ml{അനുസരണയുള്ള} (anusaraṇayuḷḷa)}{അനുസരണയുള്ള (anusaraṇayuḷḷa)}]{\ml{അനുസരണയുള്ള} (anusaraṇayuḷḷa) — obedient}

\label{lesson:ML-C206-anusaranayulla}



\begin{quote}
Before the new one: say the Malayalam for modern, then the Malayalam for naughty.
\end{quote}

\subsection*{You'll want to know: \ml{അനുസരണയുള്ള}}
\textbf{\ml{അനുസരണയുള്ള}} — \emph{anusaraṇayuḷḷa} — \textquotedblleft{}obedient\textquotedblright{}.

\begin{cousinweb}
\textbf{\ml{അനുസരണയുള്ള}} (\emph{anusaraṇayuḷḷa}) joins Sanskrit \emph{anusaraṇa}, \textquotedblleft{}following\textquotedblright{}, to native \emph{-uḷḷa}, \textquotedblleft{}that has\textquotedblright{} (DEDR 697): \textquotedblleft{}having following\textquotedblright{} is obedience.
\end{cousinweb}

\begin{grammarlens}[title={What you've built}]
One more word for this chapter.
\end{grammarlens}

\subsection*{Your turn}
\begin{itemize}
\raggedright
  \item \emph{Say it:} \emph{anusaraṇayuḷḷa}
  \item \emph{Say it:} \emph{anusaraṇayuḷḷa}, once more
  \item \emph{Say it:} say \emph{anusaraṇayuḷḷa}
  \item \emph{Recall:} say the Malayalam for modern, then the Malayalam for naughty, then say \emph{anusaraṇayuḷḷa} again
\end{itemize}

\begin{tcolorbox}[breakable,colback=teal!4,colframe=teal!35!black,title={Before you move on}]
What does \ml{അനുസരണയുള്ള} mean? (\textquotedblleft{}Obedient\textquotedblright{}.) Say it once more.
\end{tcolorbox}
\section[\texorpdfstring{\ml{സ്വാർത്ഥമായ} (svārtthamāya)}{സ്വാർത്ഥമായ (svārtthamāya)}]{\ml{സ്വാർത്ഥമായ} (svārtthamāya) — selfish}

\label{lesson:ML-C206-svartthamaya}



\begin{quote}
Before the new one: say the Malayalam for naughty, then the Malayalam for obedient.
\end{quote}

\subsection*{You'll want to know: \ml{സ്വാർത്ഥമായ}}
\textbf{\ml{സ്വാർത്ഥമായ}} — \emph{svārtthamāya} — \textquotedblleft{}selfish\textquotedblright{}.

\begin{cousinweb}
\textbf{\ml{സ്വാർത്ഥമായ}} (\emph{svārtthamāya}) is Sanskrit \emph{svārtha}, \textquotedblleft{}one's own aim\textquotedblright{}: \emph{sva}, \textquotedblleft{}own\textquotedblright{}, and \emph{artha}, \textquotedblleft{}aim\textquotedblright{} (CDIAL 638), with native \emph{āya}, \textquotedblleft{}having become\textquotedblright{} (DEDR 333). DEDR 333 compares Sanskrit \emph{ām}.
\end{cousinweb}

\begin{grammarlens}[title={What you've built}]
One more word for this chapter.
\end{grammarlens}

\subsection*{Your turn}
\begin{itemize}
\raggedright
  \item \emph{Say it:} \emph{svārtthamāya}
  \item \emph{Say it:} \emph{svārtthamāya}, once more
  \item \emph{Say it:} say \emph{svārtthamāya}
  \item \emph{Recall:} say the Malayalam for naughty, then the Malayalam for obedient, then say \emph{svārtthamāya} again
\end{itemize}

\begin{tcolorbox}[breakable,colback=teal!4,colframe=teal!35!black,title={Before you move on}]
What does \ml{സ്വാർത്ഥമായ} mean? (\textquotedblleft{}Selfish\textquotedblright{}.) Say it once more.
\end{tcolorbox}
\section[\texorpdfstring{\ml{സ്ഥിരമായ} (sthiramāya)}{സ്ഥിരമായ (sthiramāya)}]{\ml{സ്ഥിരമായ} (sthiramāya) — permanent, regular}

\label{lesson:ML-C206-sthiramaya}



\begin{quote}
Before the new one: say the Malayalam for obedient, then the Malayalam for selfish.
\end{quote}

\subsection*{You'll want to know: \ml{സ്ഥിരമായ}}
\textbf{\ml{സ്ഥിരമായ}} — \emph{sthiramāya} — \textquotedblleft{}permanent, regular\textquotedblright{}.

\begin{cousinweb}
\textbf{\ml{സ്ഥിരമായ}} (\emph{sthiramāya}) joins Sanskrit \emph{sthira}, \textquotedblleft{}firm, durable\textquotedblright{} (CDIAL 13771), to native \emph{āya}, \textquotedblleft{}having become\textquotedblright{} (DEDR 333). DEDR 333 compares Sanskrit \emph{ām}.
\end{cousinweb}

\begin{grammarlens}[title={What you've built}]
One more word for this chapter.
\end{grammarlens}

\subsection*{Your turn}
\begin{itemize}
\raggedright
  \item \emph{Say it:} \emph{sthiramāya}
  \item \emph{Say it:} \emph{sthiramāya}, once more
  \item \emph{Say it:} say \emph{sthiramāya}
  \item \emph{Recall:} say the Malayalam for obedient, then the Malayalam for selfish, then say \emph{sthiramāya} again
\end{itemize}

\begin{tcolorbox}[breakable,colback=teal!4,colframe=teal!35!black,title={Before you move on}]
What does \ml{സ്ഥിരമായ} mean? (\textquotedblleft{}Permanent, regular\textquotedblright{}.) Say it once more.
\end{tcolorbox}
\section[\texorpdfstring{\ml{ഒറ്റ} (oṟṟa)}{ഒറ്റ (oṟṟa)}]{\ml{ഒറ്റ} (oṟṟa) — single, lone}

\label{lesson:ML-C206-orra}



\begin{quote}
Before the new one: say the Malayalam for selfish, then the Malayalam for permanent.
\end{quote}

\subsection*{You'll want to know: \ml{ഒറ്റ}}
\textbf{\ml{ഒറ്റ}} — \emph{oṟṟa} — \textquotedblleft{}single, lone\textquotedblright{}.

\begin{cousinweb}
\textbf{\ml{ഒറ്റ}} (\emph{oṟṟa}) is native Dravidian: DEDR 990 lists it, \textquotedblleft{}one, single, odd\textquotedblright{}, under the word for \textquotedblleft{}one\textquotedblright{}, with Tamil \emph{oṟṟai}.
\end{cousinweb}

\begin{grammarlens}[title={What you've built}]
One more word for this chapter.
\end{grammarlens}

\subsection*{Your turn}
\begin{itemize}
\raggedright
  \item \emph{Say it:} \emph{oṟṟa}
  \item \emph{Say it:} \emph{oṟṟa}, once more
  \item \emph{Say it:} say \emph{oṟṟa}
  \item \emph{Recall:} say the Malayalam for selfish, then the Malayalam for permanent, then say \emph{oṟṟa} again
\end{itemize}

\begin{tcolorbox}[breakable,colback=teal!4,colframe=teal!35!black,title={Before you move on}]
What does \ml{ഒറ്റ} mean? (\textquotedblleft{}Single, lone\textquotedblright{}.) Say it once more.
\end{tcolorbox}
